\section{Бэкстеппинг}
\label{sec:backstepping}

\subsection{Виртуальное управление: что именно строится}
\textit{L5, слайды 1--6.}

Бэкстеппинг позволяет последовательно стабилизировать систему,
в которой управление достигает нужной переменной через другие состояния.
Основной шаг лекции рассматривает систему
\[
  \dot\eta=f(\eta)+g(\eta)\xi,\qquad \dot\xi=u,
  \qquad \eta\in\mathbb R^n,\quad \xi,u\in\mathbb R.
\]
Здесь $f,g\in\mathbb R^n$. В уравнении для $\eta$ величина $\xi$
играет роль входа, однако физически это состояние, а настоящий вход $u$
задаёт лишь скорость изменения $\xi$.

Сначала временно рассматриваем $\xi$ как свободно выбираемое
\emph{виртуальное управление}. Требуется найти гладкую функцию
$\xi=\phi(\eta)$, $\phi(0)=0$, и положительно определённую функцию
Ляпунова $V(\eta)$ так, чтобы
\[
  \frac{\partial V}{\partial\eta}
  \bigl[f(\eta)+g(\eta)\phi(\eta)\bigr]
  \le -W(\eta),
\]
где $W$ положительно определена. Это решает задачу только для
подсистемы $\eta$: состояние $\xi$ ещё не обязано совпадать с $\phi$.
Поэтому вводится ошибка
\[
  z=\xi-\phi(\eta),\qquad \xi=\phi(\eta)+z.
\]
Прибавление и вычитание $g(\eta)\phi(\eta)$ даёт тождественную запись
исходного уравнения:
\[
  \dot\eta=f(\eta)+g(\eta)\phi(\eta)+g(\eta)z.
\]
Первые два слагаемых описывают уже стабилизированную подсистему,
а последнее показывает, как ошибка реализации виртуального управления
воздействует на неё. Следующий шаг должен одновременно погасить $z$
и учесть это воздействие.

\subsection{Один шаг: цепное правило, составная функция и закон управления}
\textit{L5, слайды 7--11.}

Дифференцируя определение ошибки, получаем
\[
  \dot z=u-\dot\phi,\qquad
  \dot\phi=
  \frac{\partial\phi}{\partial\eta}\dot\eta
  =\frac{\partial\phi}{\partial\eta}
     \bigl[f(\eta)+g(\eta)\xi\bigr].
\]
Это полная производная функции $\phi(\eta(t))$ по времени.
В неё подставляется \emph{фактическое} уравнение для $\dot\eta$,
а не идеальный режим $\xi=\phi(\eta)$: ошибка $z$ может быть ненулевой.
Временная замена входа $v=u-\dot\phi$ приводит к системе
\[
  \dot\eta=f+g\phi+gz,\qquad \dot z=v.
\]
Для неё строим составную функцию Ляпунова
\[
  V_c(\eta,\xi)=V(\eta)+\frac12z^2.
\]
Член $z^2/2$ оценивает ошибку, не учитываемую исходной функцией $V$.
По цепному правилу
\begin{align*}
  \dot V_c
  &=\frac{\partial V}{\partial\eta}(f+g\phi+gz)+zv\\
  &\le -W(\eta)
       +z\left(\frac{\partial V}{\partial\eta}g+v\right).
\end{align*}
Слагаемое $z(\partial V/\partial\eta)g$ не имеет определённого знака.
Выбираем
\[
  v=-\frac{\partial V}{\partial\eta}g-kz,\qquad k>0.
\]
Первое слагаемое компенсирует перекрёстный член, второе создаёт
отрицательный квадрат. Возвращаясь к настоящему входу, получаем
\[
  \boxed{
  u=\frac{\partial\phi}{\partial\eta}(f+g\xi)
    -\frac{\partial V}{\partial\eta}g
    -k\bigl[\xi-\phi(\eta)\bigr]
  },\qquad
  \dot V_c\le -W(\eta)-kz^2.
\]
Таким образом, $\dot\phi$ заставляет $\xi$ следовать за изменяющимся
значением $\phi$, компенсация перекрёстного члена связывает два шага
синтеза, а $-kz$ обеспечивает затухание ошибки. Простая команда
$u=-kz$ в общем случае этих трёх задач не решает.

\subsection{Условия леммы и смысл глобального результата}
\textit{L5, слайд 12.}

Лемма бэкстеппинга требует стабилизирующую функцию $\phi$ с
$\phi(0)=0$, положительно определённую $V$ и оценку её производной
с положительно определённой $W$. Функции должны быть достаточно
гладкими для использованных производных и существования решений.
Если условия выполнены в некоторой окрестности, результат локальный.
Если они выполнены во всём пространстве и $V$ радиально неограниченна,
начало координат глобально асимптотически устойчиво.

\textbf{Дополнение для защиты.}
Радиальная неограниченность означает $V(\eta)\to\infty$ при
$\norm{\eta}\to\infty$. Она исключает уход состояния на бесконечность
при ограниченном значении функции Ляпунова. Замена
$(\eta,\xi)\mapsto(\eta,z)$ глобально обратима:
$\xi=z+\phi(\eta)$. Если $V_c$ ограничена, ограничены $\eta$ и $z$;
на ограниченном множестве $\eta$ ограничена и непрерывная $\phi$,
следовательно, ограничена $\xi$. Поэтому подуровни $V_c$ компактны
и в исходных координатах. Вместе с
$\dot V_c\le-W-kz^2$ это даёт глобальную асимптотическую устойчивость.

Отрицательно определённая производная доказывает асимптотическую
устойчивость; для экспоненциальной оценки нужна дополнительная связь,
например $\dot V_c\le-cV_c$, $c>0$. Даже экспоненциальная оценка ошибки
в нелинейных координатах сама по себе не задаёт глобальную линейную
оценку нормы исходного состояния.

\subsection{Первый пример лекции: второй порядок}
\textit{L5, слайды 13--17.}

Рассматривается система
\[
  \dot x_1=x_1^2-x_1^3+x_2,\qquad \dot x_2=u.
\]
На первом шаге $x_2$ считается виртуальным управлением. Выбор
\[
  \phi(x_1)=-x_1^2-x_1,\qquad V(x_1)=\frac12x_1^2
\]
даёт $\dot x_1=-x_1-x_1^3$ и
$\dot V=-x_1^2-x_1^4$. Нелинейное слагаемое $-x_1^3$ сохраняется:
оно уже уменьшает $V$, поэтому компенсировать его не требуется.

Вводим ошибку и дифференцируем её:
\begin{align*}
  z_2&=x_2-\phi(x_1)=x_2+x_1^2+x_1,\\
  \dot x_1&=-x_1-x_1^3+z_2,\\
  \dot\phi&=-(1+2x_1)(-x_1-x_1^3+z_2),\\
  \dot z_2&=u+(1+2x_1)(-x_1-x_1^3+z_2).
\end{align*}
При $V_c=\tfrac12x_1^2+\tfrac12z_2^2$ имеем
\begin{align*}
  \dot V_c
  &=x_1(-x_1-x_1^3+z_2)
    +z_2\bigl[u+(1+2x_1)(-x_1-x_1^3+z_2)\bigr]\\
  &=-x_1^2-x_1^4+
    z_2\bigl[x_1+u+(1+2x_1)(-x_1-x_1^3+z_2)\bigr].
\end{align*}
Чтобы последняя скобка стала равной $-z_2$, выбираем
\[
  u=-x_1-(1+2x_1)(-x_1-x_1^3+z_2)-z_2.
\]
В исходных координатах тот же закон имеет вид
\[
  u=-x_1-(1+2x_1)(x_1^2-x_1^3+x_2)
    -(x_2+x_1^2+x_1).
\]
Получаем
\[
  \dot V_c=-x_1^2-x_1^4-z_2^2
  =-x_1^2-x_1^4-(x_2+x_1^2+x_1)^2.
\]
Все условия леммы глобальны, замена координат обратима, а функция
Ляпунова радиально неограниченна. Поэтому начало координат глобально
асимптотически устойчиво. Последний член $-z_2$ в управлении нужен
именно для появления квадрата $-z_2^2$ в производной.

\subsection{Второй пример лекции: добавление ещё одного интегратора}
\textit{L5, слайды 18--24.}

Теперь система имеет вид
\[
  \dot x_1=x_1^2-x_1^3+x_2,\qquad
  \dot x_2=x_3,\qquad \dot x_3=u.
\]
Уже найденный регулятор второго порядка становится виртуальным
управлением $x_3$:
\[
  x_3=\phi(x_1,x_2)
  =-x_1-(1+2x_1)(x_1^2-x_1^3+x_2)-(x_2+x_1^2+x_1).
\]
Здесь обозначение $\phi$ относится к \emph{новому} шагу: теперь оно
зависит от двух состояний. Функция Ляпунова стабилизированной
двумерной подсистемы равна
\[
  V(x_1,x_2)=\frac12x_1^2+\frac12(x_2+x_1^2+x_1)^2.
\]
Обозначая $z_2=x_2+x_1^2+x_1$, вводим следующую ошибку
\[
  z_3=x_3-\phi(x_1,x_2),\qquad
  \dot x_2=\phi+z_3,
\]
\[
  \dot z_3=u-\dot\phi,
  \qquad
  \dot\phi=
  \frac{\partial\phi}{\partial x_1}(x_1^2-x_1^3+x_2)
  +\frac{\partial\phi}{\partial x_2}(\phi+z_3).
\]
Второе слагаемое цепного правила обязательно: $\phi$ зависит и от
$x_2$, а фактическое $\dot x_2$ равно $x_3=\phi+z_3$.
При $V_c=V+\tfrac12z_3^2$ производная равна
\begin{align*}
  \dot V_c
  &=\frac{\partial V}{\partial x_1}\dot x_1
    +\frac{\partial V}{\partial x_2}\dot x_2
    +z_3(u-\dot\phi)\\
  &=-x_1^2-x_1^4-z_2^2
    +z_3\left(\frac{\partial V}{\partial x_2}+u-\dot\phi\right).
\end{align*}
Поскольку $\partial V/\partial x_2=z_2$, лекционная формула даёт
\begin{align*}
  u&=-\frac{\partial V}{\partial x_2}+\dot\phi-z_3\\
   &=-z_2+
     \frac{\partial\phi}{\partial x_1}(x_1^2-x_1^3+x_2)
     +\frac{\partial\phi}{\partial x_2}(\phi+z_3)-z_3.
\end{align*}
Итак,
\[
  \dot V_c=-x_1^2-x_1^4-(x_2+x_1^2+x_1)^2-z_3^2.
\]
Отрицательные слагаемые предыдущего шага сохранены, а новый шаг
добавил $-z_3^2$. В этом и состоит рекурсивность метода: готовые
виртуальное управление и функция Ляпунова используются как исходные
данные для следующего интегратора.
Обе замены ошибок глобально обратимы, а все функции гладкие;
поэтому и для этой системы начало координат глобально
асимптотически устойчиво.

\newpage
\subsection{Общий вид strict-feedback и ненулевые коэффициенты}
\textit{L5, слайды 25--27.}

В обобщении лекции рассматривается каскад
\begin{align*}
  \dot x&=f_0(x)+g_0(x)z_1,\\
  \dot z_1&=f_1(x,z_1)+g_1(x,z_1)z_2,\\
  \dot z_2&=f_2(x,z_1,z_2)+g_2(x,z_1,z_2)z_3,\\
  &\ \vdots\\
  \dot z_{k-1}&=f_{k-1}(x,z_1,\ldots,z_{k-1})
     +g_{k-1}(x,z_1,\ldots,z_{k-1})z_k,\\
  \dot z_k&=f_k(x,z_1,\ldots,z_k)
     +g_k(x,z_1,\ldots,z_k)u,
\end{align*}
где $g_i\ne0$ на области синтеза для $i=1,\ldots,k$.
Каждое уравнение зависит от предшествующих состояний и линейно от
следующего состояния; последнее уравнение линейно зависит от $u$.
На этом слайде $z_i$ обозначают исходные координаты каскада,
а не уже построенные ошибки. При применении метода разность
между очередным состоянием и виртуальным управлением вводится отдельно.

\textbf{Дополнение для защиты.}
Если очередной шаг имеет вид
\[
  \dot\eta=f(\eta)+g(\eta)\xi,\qquad
  \dot\xi=f_i(\eta,\xi)+g_i(\eta,\xi)u,
\]
то для $z=\xi-\phi(\eta)$ получаем
$\dot z=f_i+g_i u-\dot\phi$. Поэтому обобщённая формула равна
\[
  u=\frac{\dot\phi-f_i-
       (\partial V/\partial\eta)g-kz}{g_i}.
\]
Деление на $g_i$ объясняет условие $g_i\ne0$.
Если коэффициент обращается в нуль, этот закон там не определён,
и глобальный результат из него не следует.
Знак $g_i$ может быть отрицательным: его учитывают при компенсации
перекрёстных членов.

Слайд 26 показывает применение адаптивного бэкстеппинга к квадрокоптеру:
эксперимент с вентилятором сопровождается графиком нормы ошибки.
Это иллюстрация применения метода; модель квадрокоптера и уравнения
адаптации на слайде не выводятся. Слайд 27 завершает лекцию.

\subsection{Лабораторная 4, система 1: как пройти два шага}
\textit{L5, слайды 3--12 и 25. Дополнение для защиты: разбор Work4.}

Исходная система
\[
  \dot x_1=x_2+\sin x_1+x_1^2,\qquad
  \dot x_2=x_1^2+g_1(x_1)u,\qquad g_1=2+\sin x_1.
\]
На первом шаге хотим получить $\dot x_1=-x_1$. Для этого при
виртуальном выборе $x_2=\phi(x_1)$ должны компенсироваться
$\sin x_1+x_1^2$ и появиться $-x_1$:
\[
  \phi(x_1)=-x_1-\sin x_1-x_1^2,\qquad
  V=\frac12x_1^2,\qquad \dot V=-x_1^2.
\]
Настоящее $x_2$ отличается от $\phi$, поэтому
\[
  z_2=x_2-\phi,\qquad
  \dot x_1=-x_1+z_2,\qquad
  \dot z_2=x_1^2+g_1u-\dot\phi.
\]
Производную виртуального управления вычисляем до подстановки
идеального режима:
\[
  \dot\phi=(-1-\cos x_1-2x_1)
             (x_2+\sin x_1+x_1^2).
\]
Функция $V_c=\tfrac12(x_1^2+z_2^2)$ даёт
\[
  \dot V_c=-x_1^2+
     z_2\bigl[x_1+x_1^2+g_1u-\dot\phi\bigr].
\]
Чтобы скобка стала $-z_2$, используем
\[
  \boxed{u=\frac{\dot\phi-x_1^2-x_1-z_2}{g_1}}.
\]
Слагаемое $-x_1^2$ компенсирует дрейф второго уравнения,
$\dot\phi$ учитывает движение виртуальной цели,
$-x_1$ компенсирует перекрёстный член, а $-z_2$ гасит ошибку.
После подстановки
\[
  \dot x_1=-x_1+z_2,\qquad
  \dot z_2=-x_1-z_2,\qquad
  \dot V_c=-x_1^2-z_2^2.
\]
При всех $x_1$ выполняется $g_1\ge1$; деление безопасно на всём
пространстве. Обратная замена $x_2=z_2+\phi(x_1)$ глобальна,
поэтому доказана глобальная асимптотическая устойчивость исходного нуля.

\subsection{Лабораторная 4, система 2: полезную нелинейность сохраняем}
\textit{L5, слайды 13--17. Дополнение для защиты: разбор Work4.}

Для системы
\[
  \dot x_1=x_2-x_1^3,\qquad \dot x_2=x_1+u
\]
выбираем $\phi(x_1)=-x_1$. Тогда при $x_2=\phi$
\[
  \dot x_1=-x_1-x_1^3,\qquad
  V=\frac12x_1^2,\qquad \dot V=-x_1^2-x_1^4.
\]
Как и в первом примере лекции, слагаемое $-x_1^3$ сохраняем:
оно создаёт дополнительное отрицательное $-x_1^4$.
При $z_2=x_2+x_1$ получаем
\begin{align*}
  \dot x_1&=-x_1-x_1^3+z_2,\\
  \dot\phi&=-\dot x_1,\\
  \dot z_2&=x_1+u+\dot x_1=u-x_1^3+z_2.
\end{align*}
Следовательно,
\[
  V_c=\frac12(x_1^2+z_2^2),\qquad
  \dot V_c=-x_1^2-x_1^4+
          z_2(x_1+u-x_1^3+z_2).
\]
Выбираем скобку равной $-z_2$:
\[
  \boxed{u=x_1^3-x_1-2z_2=x_1^3-3x_1-2x_2}.
\]
Тогда
\[
  \dot z_2=-x_1-z_2,\qquad
  \dot V_c=-x_1^2-x_1^4-z_2^2.
\]
Хотя в управлении присутствует положительное $x_1^3$, замкнутая
система устойчива: знак отдельного члена управления не определяет
устойчивость. Подстановка во все уравнения и вычисление $\dot V_c$
учитывают его действительную роль. Здесь условия глобальности выполнены.

\subsection{Лабораторная 4, система 3: равновесие и четыре шага}
\textit{L5, слайды 7--12, 18--25. Дополнение для защиты: разбор Work4.}

Рассматривается система
\[
  \begin{cases}
    \dot x_1=\cos x_1-x_2,\\
    \dot x_2=x_1+x_3,\\
    \dot x_3=x_1x_3+g_2(x_3)x_4,\qquad g_2=2-\sin x_3,\\
    \dot x_4=x_2x_3+2u.
  \end{cases}
\]
Сначала проверяем целевую точку. При $\bm x=0$ имеем
$\dot x_1=1$ независимо от $u$, поэтому никакое управление не сделает
исходный нуль равновесием. Полагая $x_1^*=0$, из последовательных
уравнений равновесия получаем
\[
  x_2^*=1,\qquad x_3^*=0,\qquad x_4^*=0,\qquad u^*=0.
\]
Стабилизируется точка $\bm x^*=(0,1,0,0)^T$.
Лемма формулируется для нуля, поэтому можно мысленно перейти к
отклонению $\widetilde x_2=x_2-1$. Тогда первое виртуальное управление
равно $\widetilde\phi_1=\cos x_1-1+x_1$ и обращается в нуль при
$x_1=0$. Ниже записан эквивалентный синтез сразу в исходных координатах.

\textbf{Первый шаг: выбираем $x_2$.}
При $V_1=x_1^2/2$ хотим получить $\dot x_1=-x_1$. Из-за знака
минус перед $x_2$ выбираем
\[
  \phi_1=\cos x_1+x_1,\qquad z_2=x_2-\phi_1.
\]
Отсюда
\[
  \dot x_1=-x_1-z_2,\qquad
  \dot V_1=-x_1^2-x_1z_2,\qquad
  \dot\phi_1=(1-\sin x_1)(\cos x_1-x_2).
\]
Знак перекрёстного члена теперь отрицательный; это повлияет
на требуемый знак связи в следующем уравнении.

\textbf{Второй шаг: выбираем $x_3$.}
Для $V_2=V_1+z_2^2/2$ имеем
\begin{align*}
  \dot z_2&=x_1+x_3-\dot\phi_1,\\
  \dot V_2&=-x_1^2-x_1z_2
          +z_2(x_1+x_3-\dot\phi_1)
           =-x_1^2+z_2(x_3-\dot\phi_1).
\end{align*}
Поэтому выбираем
\[
  \phi_2=\dot\phi_1-z_2,\qquad z_3=x_3-\phi_2.
\]
Тогда
\[
  \dot z_2=x_1-z_2+z_3,\qquad
  \dot V_2=-x_1^2-z_2^2+z_2z_3.
\]
Член $+x_1$ в $\dot z_2$ компенсирует $-x_1z_2$
из производной $V_1$.

\textbf{Третий шаг: выбираем $x_4$.}
Вводим $V_3=V_2+z_3^2/2$. Поскольку
$\dot z_3=x_1x_3+g_2x_4-\dot\phi_2$, получаем
\[
  \dot V_3=-x_1^2-z_2^2+
    z_3\bigl[z_2+x_1x_3+g_2x_4-\dot\phi_2\bigr].
\]
Виртуальное управление и очередная ошибка равны
\[
  \phi_3=\frac{\dot\phi_2-x_1x_3-z_2-z_3}{g_2},
  \qquad z_4=x_4-\phi_3.
\]
Отсюда
\[
  \dot z_3=-z_2-z_3+g_2z_4,\qquad
  \dot V_3=-x_1^2-z_2^2-z_3^2+g_2z_3z_4.
\]
Коэффициент $g_2$ перед последней ошибкой сохраняется; заменять его
единицей нельзя.

\textbf{Четвёртый шаг: выбираем настоящее управление $u$.}
Для $V_c=V_3+z_4^2/2$ имеем
\[
  \dot z_4=x_2x_3+2u-\dot\phi_3,
\]
\[
  \dot V_c=-x_1^2-z_2^2-z_3^2+
     z_4\bigl[g_2z_3+x_2x_3+2u-\dot\phi_3\bigr].
\]
Выбираем
\[
  \boxed{u=\frac12
    \bigl(\dot\phi_3-x_2x_3-g_2z_3-z_4\bigr)}.
\]
Тогда $\dot z_4=-g_2z_3-z_4$ и
\[
  V_c=\frac12(x_1^2+z_2^2+z_3^2+z_4^2),\qquad
  \dot V_c=-x_1^2-z_2^2-z_3^2-z_4^2.
\]
Перекрёстные члены сокращаются попарно:
$-x_1z_2+x_1z_2$, $z_2z_3-z_2z_3$ и
$g_2z_3z_4-g_2z_3z_4$.

\subsection{Как вычислить производные третьего регулятора}
\textit{L5, слайды 7, 11, 21--24. Дополнение для защиты.}

Обозначение $\dot\phi_i$ не требует измерять производные состояний.
Это функция измеренных $x_j$, полученная по цепному правилу и модели.
Поскольку $\phi_2$ зависит только от $(x_1,x_2)$,
а $\phi_3$ только от $(x_1,x_2,x_3)$,
\begin{align*}
  \dot\phi_2&=
     \frac{\partial\phi_2}{\partial x_1}(\cos x_1-x_2)
     +\frac{\partial\phi_2}{\partial x_2}(x_1+x_3),\\
  \dot\phi_3&=
     \frac{\partial\phi_3}{\partial x_1}(\cos x_1-x_2)
     +\frac{\partial\phi_3}{\partial x_2}(x_1+x_3)
     +\frac{\partial\phi_3}{\partial x_3}(x_1x_3+g_2x_4).
\end{align*}
Настоящий вход $u$ появляется только в четвёртом уравнении;
оно здесь не используется. Поэтому формула для $u$ явная,
без скрытого уравнения относительно самого $u$.

Для вычисления градиентов удобно сохранить промежуточные выражения,
не раскрывая все произведения. Введём обозначения
\[
  a=1-\sin x_1,\qquad f_1=\cos x_1-x_2,
  \qquad f_2=x_1+x_3.
\]
Тогда $\phi_2=af_1-z_2$, и её частные производные равны
\[
  p=\frac{\partial\phi_2}{\partial x_1}
    =a^2-\cos x_1\,f_1,\qquad
  q=\frac{\partial\phi_2}{\partial x_2}=-a-1.
\]
Здесь при дифференцировании $z_2=x_2-\cos x_1-x_1$ учтены
$\partial z_2/\partial x_1=-a$ и
$\partial z_2/\partial x_2=1$. Следовательно,
\[
  \dot\phi_2=pf_1+qf_2,\qquad
  N=pf_1+qf_2-x_1x_3-z_2-z_3,\qquad \phi_3=\frac{N}{g_2}.
\]
Для полного вычисления градиента достаточно следующих коротких формул:
\begin{align*}
  p_{x_1}&=-2a\cos x_1+\sin x_1\,f_1+\cos x_1\sin x_1,\\
  N_{x_1}&=p_{x_1}f_1-p\sin x_1+\cos x_1\,f_2+q-x_3+a+p,\\
  N_{x_2}&=\cos x_1\,f_1-p-1+q,\\
  N_{x_3}&=q-x_1-1.
\end{align*}
Поскольку $g_2$ зависит только от $x_3$,
\[
  \frac{\partial\phi_3}{\partial x_1}=\frac{N_{x_1}}{g_2},
  \qquad
  \frac{\partial\phi_3}{\partial x_2}=\frac{N_{x_2}}{g_2},
  \qquad
  \frac{\partial\phi_3}{\partial x_3}
     =\frac{N_{x_3}}{g_2}+\frac{N\cos x_3}{g_2^2}.
\]
Плюс во втором слагаемом последней формулы возникает из двух минусов:
$(1/g_2)'=-g_2'/g_2^2$, а $g_2'=-\cos x_3$.
Если не продифференцировать знаменатель, рассчитанное управление
перестанет обеспечивать требуемую производную $V_c$.

\subsection{Глобальность и чтение результатов моделирования лабораторной 4}
\textit{L5, слайд 12 и общий вид на слайде 25.
Дополнение для защиты.}

В третьей системе $1\le g_2\le3$, а коэффициент при $u$ равен двум.
Поэтому деления в $\phi_3$ и $u$ не создают особых точек.
Замена $(x_1,x_2,x_3,x_4)\mapsto(x_1,z_2,z_3,z_4)$ треугольна;
её обратное преобразование восстанавливает состояния последовательно:
\[
  x_2=z_2+\phi_1(x_1),\qquad
  x_3=z_3+\phi_2(x_1,x_2),\qquad
  x_4=z_4+\phi_3(x_1,x_2,x_3).
\]
Оно гладко и определено всюду. При нулевых ошибках получаем
$x_1=0$, $x_2=1$, $x_3=x_4=0$.
Ограниченность $V_c$ последовательно ограничивает $x_1,z_2,z_3,z_4$
и затем все исходные состояния. Следовательно, $V_c$ радиально
неограниченна относительно $\bm x^*$ и доказательство глобально.

В лабораторной моделирование выполнялось методом Рунге-Кутты
четвёртого порядка с шагом $0{,}001$ с на интервале $[0,10]$ с.
Для систем 1 и 2 выбрано $\bm x(0)=(1;-0{,}5)^T$,
для системы 3 --- $\bm x(0)=(0{,}5;1;-0{,}5;0{,}25)^T$.
В первых двух случаях на графиках показаны состояния,
в третьем --- отклонения $(x_1,x_2-1,x_3,x_4)$.
Их стремление к нулю означает $x_2\to1$ в третьей системе,
а не $x_2\to0$. Во всех трёх случаях управление стремится к нулю.

Для систем 1 и 3 выполнено $\dot V_c=-2V_c$, поэтому
$V_c(t)=V_c(0)e^{-2t}$; для системы 2 имеем
$\dot V_c=-2V_c-x_1^4\le-2V_c$.
Отсюда следует затухание нормы соответствующего вектора ошибок
по оценке $\norm{\bm z(t)}\le e^{-t}\norm{\bm z(0)}$,
где $\bm z=(x_1,z_2)^T$ либо $(x_1,z_2,z_3,z_4)^T$.
Отдельное исходное состояние при этом может
временно возрастать: доказательство относится к составной функции,
а не к монотонности каждой координаты.
Графики показывают один выбранный опыт; глобальный вывод обеспечивают
функция Ляпунова, отрицательная производная и глобальность замены.
